\begin{helper}
Fix natural numbers $k$ and $D$, and let $(\mathcal P,p)$ be a bounded
local pattern from \noderef{KUniformKSimpleBoundedLocalPattern}.  For any
ambient universe levels, there is a hypergraph-edge pair
$(\mathcal F,f)$ in those universe levels, still belonging to $H(k,k,D)$,
such that
\[
b_{L(\mathcal P),kD}(p)=b_{L(\mathcal F),kD}(f).
\]
\end{helper}
\begin{proof}
The bounded pattern \noderef{KUniformKSimpleBoundedLocalPattern} is a
finite labelled hypergraph on the fixed finite edge-label set
$\{0,\ldots,1+k(D-1)-1\}$ and the fixed finite vertex-label set
$\{0,\ldots,k+k^2(D-1)-1\}$, together with the distinguished edge label
$p$.  To place this data in arbitrary ambient universe levels, replace each
edge label and each vertex label by its canonical lifted copy in the desired
universe.  This gives finite edge and vertex types in those universes, with
decidable equality, and defines a hypergraph $\mathcal F$ by transporting
the edge assignment of $\mathcal P$ along the two lift equivalences.  Let
$f$ be the lifted copy of $p$.

Because the construction is just relabeling by bijections, it preserves all
incidence relations.  Therefore every edge of $\mathcal F$ has the same
cardinality as the corresponding edge of $\mathcal P$, every intersection
of two edges in $\mathcal F$ has the same cardinality as the corresponding
intersection in $\mathcal P$, and every vertex degree in $\mathcal F$ is
the same as the corresponding vertex degree in $\mathcal P$.  Since
\noderef{KUniformKSimpleBoundedLocalPattern} includes the assertion that
$\mathcal P\in H(k,k,D)$, the three components of
\noderef{HypergraphClass}, interpreted through \noderef{UniformHypergraph},
\noderef{TSimpleHypergraph}, \noderef{MaxDegreeAtMost}, and
\noderef{HypergraphDegree}, transport to prove
$\mathcal F\in H(k,k,D)$.

The same relabeling gives an isomorphism between the line graphs
$L(\mathcal P)$ and $L(\mathcal F)$ around the distinguished labels:
by \noderef{LineGraphOfHypergraph}, adjacency is exactly nonempty
intersection of the corresponding hyperedges, and nonempty intersection is
preserved and reflected by the vertex-label bijection.  Hence the
line-graph degree of $p$ equals that of $f$, the independent pairs in the
neighborhood counted by \noderef{IndependentPairCount} correspond
bijectively, and the independent triples counted by
\noderef{IndependentTripleCount} correspond bijectively.  The formula
\noderef{LocalBParameter} depends only on these three counts and the common
parameter $kD$, so the two displayed local $b$-parameters are equal.
\end{proof}
